\documentclass[11pt,a4paper]{article}
\usepackage[margin=2.6cm]{geometry}
\usepackage{amsmath,amssymb}
\usepackage{booktabs}
\usepackage{enumitem}
\usepackage{xcolor}
\usepackage[colorlinks=true,linkcolor=teal!70!black,citecolor=teal!70!black,urlcolor=teal!70!black]{hyperref}

% Methodology note (session 15, 2026-10-29): the "How this work was done" section of counting-full-v2.tex, moved out of the physics papers and expanded.
% Cited by toolkit-1p1-v1.tex and embedding-3p1-v1.tex as [MethodNote]. Lives in the public repository; arXiv-able (cs.CY, cross-list cs.AI; needs an endorser) if Roberto wants.

\title{Two causal-set papers written with an AI research partner:\\ methodology note}
\author{Roberto Bada\\ \small SapereIA, Barcelona, Spain\\ \small Institute of Cosmos Sciences, University of Barcelona (ICCUB), Barcelona, Spain\\ \small \texttt{roberto@sapereia.com}}
\date{Version 1, 29 October 2026}

\begin{document}
\maketitle

\begin{abstract}
This note describes how two papers on the kinematics of causal sets~\cite{Toolkit,Embedding} were produced in six days by one physicist, new to the field, working with an AI system as a research partner. We describe the division of labor, the research notebook that carried the work across sessions, the verification rules that stood in for slower human checks, the errors that the process made and caught, and the limits of what a single project of this kind can show. The full notebook, including the mistake log and the abandoned lines, is public. The note is meant to let readers judge the physics papers' reliability and to let others reuse or criticize the method; the physics itself is in the papers.
\end{abstract}

\section{Purpose}
The author works in gravitational-wave research with the Institute of Cosmos Sciences of the University of Barcelona (ICCUB) and is a member of the LIGO--Virgo--KAGRA Collaboration; neither the papers nor this note are products of that collaboration. As a co-founder of SapereIA he studies how AI systems can help advance scientific knowledge. The project was set up partly to test one question: can a physicist, new to a field, and an AI system working as a colleague reach results that are new, checked and publishable, in days rather than months? Causal set theory was chosen because it was outside the author's expertise, has a compact literature, and admits numerical tests on a laptop and a small cluster.

The physics papers state the use of AI in their acknowledgments and cite this note; they do not discuss it further, because a methods discussion of this kind is out of place in a physics journal article. This note is where it belongs.

\section{Setup}
\paragraph{System.} The AI system was Claude Opus~5.5 (Anthropic), used through the Claude Code command-line tool, which can read and write files, run code locally, and prepare jobs for a computing cluster. Cluster jobs were submitted and retrieved by the author. The AI system is not an author of the papers.

\paragraph{Sessions and notebook.} The work ran over fifteen sessions of a few hours each between 24 and 29 October 2026, as a conversation recorded in a version-controlled research notebook. The notebook has a fixed structure, read at the start of every session:
\begin{itemize}[itemsep=1pt]
\item an index of ideas with a status each (seed, exploring, developed, parked, known, killed) and a list of open questions;
\item one file per idea, in which derivations, numbers and objections were written as soon as they appeared;
\item a journal entry per session (what was done, decisions, mistakes, next steps);
\item a \emph{mistake log}: each error, and the working rule adopted to prevent it (41 entries at the time of writing; Table~\ref{tab:mistakes});
\item a \emph{graveyard}: killed ideas and dropped approaches, each with the reason; and a list of claims given back after they were found in the literature;
\item a literature file, with why each reference matters, and a notation file.
\end{itemize}
Every claim in the notebook carries a label: \emph{established} (mainstream, experimentally backed), \emph{literature} (published, not necessarily accepted), or \emph{ours} (our own construction or speculation). The project ran under written constraints set by the author (for example, no background structures and local Lorentz invariance) that every proposed model had to respect.

\paragraph{Output.} About 6500 lines of simulation code in 71 scripts, about 30 cluster jobs, ten internal drafts before the final split, each with a written review, and the two papers that this note accompanies. Table~\ref{tab:timeline} lists when the results that entered the papers appeared.

\begin{table}[h]
\centering\small
\begin{tabular}{@{}llp{0.7\linewidth}@{}}
\toprule
Day & Sessions & Results that entered the papers \\
\midrule
24 Oct & 1 & starting idea: discrete spacetime from counting \\
25 Oct & 2--6 & rectangle identity and depth drift~\cite{Toolkit}; first embedding test; first cluster runs \\
26 Oct & 7 & fixed-$n$ propagator target~\cite{Toolkit} \\
27 Oct & 8--11 & phases at $N=400$; literature crawl; null-pair fraction; $3+1$ estimator; reproduction of Johnston 2025~\cite{Embedding} \\
28 Oct & 12--13 & large-$n$ cluster runs; low-count diagnosis; likelihood fit, ablation, second shape~\cite{Embedding} \\
29 Oct & 14--15 & second shape at equal density; final literature checks; public release; split into two papers; referee-triggered runs (oracle on the benchmark's own detours, per-pair and per-element errors, $1+1$ baselines, propagator vs ordering fraction, causal fidelity) \\
\bottomrule
\end{tabular}
\caption{When each result appeared. A session is one sitting of a few hours.}\label{tab:timeline}
\end{table}

\section{Roles}
The author brought in the paper that started the distance work (Bogu\~n\'a and Krioukov's causal-overlap method), proposed the questions and several of the ideas that ended up working, set the physical constraints, and decided which lines to pursue or drop. Two examples: the suggestion to use each distance estimate only where it is accurate (a ``switch'', which was the main $3+1$ result until the likelihood fit superseded it), and the questions that led to the low-count diagnosis and to the test in a second region~\cite{Embedding}. The AI system searched and read the literature, carried out the derivations, wrote and ran the simulation code (on a laptop and, through job scripts, on the ICCUB cluster), and drafted the text. The author followed and checked the derivations, and read and approved the final texts; he did not review the code line by line. The numerical claims therefore rest on the checks of Sec.~\ref{sec:checks} and on the public code, not on independent human review of the code.

\section{Verification rules}\label{sec:checks}
Speed is worth something only if the results hold. The working rules, most of them written into the mistake log after a specific failure, fall into five groups.
\begin{enumerate}[leftmargin=*,itemsep=2pt]
\item \emph{Known answers first.} Every tool is run on a case with a known answer (flat space, de Sitter, the continuum limit) before anything else. The published numbers of a method we compare against are reproduced before any comparison is made.
\item \emph{Statistics.} Errors are run-to-run spreads over independent seeds, not per-event errors from correlated samples. Any parameter chosen while looking at data is re-tested on fresh seeds; every comparison is paired (same universes for every method); multi-term methods are ablated before the gain is attributed to a term; a perfect fit is checked for whether there was anything to fit.
\item \emph{Geometry and limits.} A first result is rerun in a second geometry before it is interpreted, and the opposite limits (no coupling, zero size, infinite size) are checked before a mechanism is claimed. Samplers are checked (acceptance rates) before their output is called physics.
\item \emph{Literature.} Citations are verified; nothing unverified is recorded without a mark. Before a claim enters a draft as new, the full citation lists of the closest papers are read and the citing papers' full texts searched; ``not found in a brief search'' is not reported as ``new''. A paper is quoted before we say what it does not do.
\item \emph{Computation.} Jobs are sized for the slowest loaded node and checkpoint their results; a short test precedes every new code on the cluster; the whole import tree is shipped with each job.
\end{enumerate}

\section{Errors made and caught}
Errors were frequent, and most were the AI system's. Table~\ref{tab:mistakes} groups the 41 entries of the mistake log. Examples that changed the papers:
\begin{itemize}[itemsep=1pt]
\item a ``hysteresis'' at the phase transition that came from Monte Carlo chains with acceptance below $10^{-3}$, i.e., a frozen sampler (withdrawn; kept in an appendix of~\cite{Toolkit});
\item a detector believed to be new that turned out to be Clough and Evans's (now credited; our contribution narrowed to its calibration and a missing statistic);
\item a comparison that had put two different distances side by side (redone on the same quantity);
\item a likelihood fit first presented as resting on a new term of ours, until removing each term in turn showed that the gain comes from the related-pair term, which is Casasola's, and that our term makes it worse (now credited, our term moved to the appendix of~\cite{Embedding});
\item a statement that Bogu\~n\'a and Krioukov had only an approximate inversion, when their exact one was in their Eq.~(22) (corrected; our claim narrowed to unequal depths).
\end{itemize}
In each case the error was found by one of the rules above, by the author's questions, or by a literature pass, and was recorded with the rule it produced.

\begin{table}[h]
\centering\small
\begin{tabular}{@{}p{0.27\linewidth}cp{0.55\linewidth}@{}}
\toprule
Group & Entries & Typical example \\
\midrule
Claims beyond the computation & 13 & an analogy stretched past its calculation; a mechanism asserted without checking a limit; a result that was the shape of the box \\
Statistics and numerical method & 11 & correlated samples or bins read as significant; ratios of counts read as unbiased; a frozen sampler read as physics; a gain attributed before ablation \\
Code and computing & 8 & a lookup table that did not cover the use case; cluster jobs sized on a quiet node; a missing import on the cluster \\
Literature and credit & 5 & a guessed identifier; novelty claimed after a shallow search; a paper's content misstated from its abstract \\
Process and communication & 4 & a draft that followed chronology instead of a storyline; an option offered that contradicted the draft's own promise \\
\bottomrule
\end{tabular}
\caption{Entries of the mistake log (39 numbered rules and two sub-rules), grouped by kind. Classification by the author of this note.}\label{tab:mistakes}
\end{table}

\section{Writing and review}
Drafts were written to tell one storyline, built from what advanced the work rather than in chronological order; abandoned lines go to an appendix. Each draft received a written review against an explicit bar (a credible arXiv preprint; later, the standard of Physical Review D), listing what a referee would push on and what was missing. The final step split the combined draft into the two papers. The AI system then alternated, over four rounds, between referee passes against the norms of the target journal and author revisions. The referees were separate instances of the same model; they saw the drafts, the data files and a prompt written by the authoring instance, which pointed them to the files and checks, but not the author's reasoning; the review was therefore independent of the authoring context but not blind. Several referee points required new runs (an oracle test run on the benchmark method itself rather than on our own estimator, per-pair and per-element errors, baselines for the $1+1$ distance estimator, a comparison of the propagator target with the ordering fraction), and claims were weakened or corrected in both papers as a result: the propagator target turned out to be no stronger than the ordering fraction as a diagnostic, a headline percentage of the $3+1$ paper moved when the oracle was rerun on matched universes, and a caveat on position errors was added. The referee reports, the final reviews and all earlier drafts are in the public notebook.

\section{What this does and does not show}
One project is an anecdote, not a measurement: there is no control, and the speed depends on the author, the field and the problem. The papers have not been refereed, and a reader who wants to verify them should run the code rather than trust the process. What the project shows is narrower: with an AI system doing the reading, the algebra and the programming, under explicit rules against self-deception, a physicist from a neighboring field went from a first idea to a reproduction of a 2025 result, a diagnosis of the source of its residual error and a substantial improvement on it~\cite{Embedding}, and to a set of exact $1+1$ tools~\cite{Toolkit}, in a few days, provided that the literature search was redone before every claim of novelty. The main risks we saw are the ones the rules target: confident claims beyond the computation, and novelty claims that a deeper search overturns.

\section*{Availability}
The code and the full notebook (idea files, session journals, mistake log, graveyard, internal drafts and their reviews) are public at \url{https://github.com/presindev/causet-counting}, as a snapshot; only notes on the computing cluster's internal configuration are omitted.

\section*{Acknowledgements}
This work was funded by SapereIA. Simulations were run on the Nyx cluster of ICCUB. The AI system drafted this note under the author's direction; the author takes full responsibility for its content.

\begin{thebibliography}{9}\small
\bibitem{Toolkit} R.~Bada, ``Counting identities in two-dimensional causal sets: spatial distance, depth-drift curvature, and a fixed-size propagator target'' (2026).
\bibitem{Embedding} R.~Bada, ``Spatial embedding of $3+1$ causal sets: low-count spacelike pairs and a likelihood fit to interval counts'' (2026).
\end{thebibliography}
\end{document}
